% Informacja o wykorzystaniu narzędzi AI — wzór z załącznika nr 2A (wariant 1), etap 1
% Kompilacja: pdflatex informacja-2a.tex
\documentclass[11pt,a4paper]{article}

\usepackage[utf8]{inputenc}
\usepackage[T1]{fontenc}
\usepackage[polish]{babel}
\usepackage[margin=2.3cm]{geometry}
\usepackage{enumitem}
\usepackage{hyperref}
\hypersetup{colorlinks=true, linkcolor=black, urlcolor=blue}
\setlist{nosep, leftmargin=1.5em}
\pagestyle{empty}
\setlength{\parindent}{0pt}
\setlength{\parskip}{0.5em}

\begin{document}

\begin{flushright}\small
Załącznik nr 2A do Zasad korzystania z narzędzi AI\\
przez studentów i doktorantów Politechniki Opolskiej
\end{flushright}

\begin{center}
{\large\bfseries INFORMACJA O WYKORZYSTANIU NARZĘDZI AI}\\[0.4em]
Wariant 1 --- korzystano z AI w zakresie wymagającym ujawnienia
\end{center}

\textbf{Przedmiot:} Rozpoznawanie wzorców (studia II stopnia, 2026/2027, grupa 2).\\
\textbf{Materiał:} etap 1 projektu --- prezentacja tematu 06 \emph{Obciążenie i zmienność oceny przez
walidację krzyżową} (slajdy, protokół badania, kod próby uruchomienia w repozytorium
\url{https://github.com/FylyppS/rozpoznawanie-wzorcow}).\\
\textbf{Poziom korzystania z AI} ustalony w kontrakcie przedmiotu: poziom 2 --- AI jako wsparcie pracy.\\
\textbf{Autorzy:} Dariusz Salzburg (103931), Filip Sejud (103928).

\medskip
W przygotowaniu \textbf{materiałów etapu 1 projektu}
korzystaliśmy z następujących narzędzi AI:

\textbf{Nazwa narzędzia lub narzędzi:}\\
Claude (Anthropic) --- asystent generatywny w środowisku programistycznym Claude Code.

\textbf{Cel i zakres wykorzystania:}
\begin{itemize}
  \item analiza publikacji R.~Kohaviego (IJCAI 1995): streszczenie metodologii i wyników;
  \item przygotowanie szkicu protokołu badania (generator danych o znanym błędzie Bayesa, schemat
    ziaren losowych, liczba powtórzeń, budżet obliczeniowy) --- decyzje o przyjęciu tych ustaleń
    oraz hipotezy badawcze są nasze;
  \item napisanie kodu eksperymentu w języku Python (pakiet \texttt{cvbias}, skrypty
    \texttt{run\_experiment.py} i \texttt{make\_figures.py}) oraz wykresów ilustrujących generator danych.
\end{itemize}

\textbf{Elementy przygotowane lub zmienione z wykorzystaniem AI:}
\begin{itemize}
  \item streszczenie i analiza publikacji Kohaviego;
  \item szkic protokołu badania;
  \item kod źródłowy w repozytorium (generator danych, przebieg eksperymentu, podsumowania) i rysunki
    generatora --- przez nas przeczytane, uruchomione i sprawdzone (m.in.\ ręczne przeliczenie średniej
    błędów foldów i błędu Bayesa).
\end{itemize}

\textbf{Bez udziału AI} przygotowaliśmy prezentację (treść i układ slajdów) oraz opracowaliśmy
zrozumienie tematu. Pomiary (uruchomienie próby, \texttt{results/pilot}) wykonaliśmy sami na naszym
komputerze i sami przeanalizowaliśmy ich wyniki; żadne wartości liczbowe nie zostały wygenerowane przez AI.

Wyniki uzyskane przy użyciu AI zostały przez nas sprawdzone, a za ostateczną treść i poprawność
przygotowanego materiału ponosimy odpowiedzialność.

\vspace{2.2em}
\noindent
\begin{tabular}{@{}p{0.48\textwidth}p{0.48\textwidth}@{}}
\dotfill & \dotfill \\
Dariusz Salzburg & Filip Sejud \\
\end{tabular}

\vspace{1em}
{\small Opole, 29 października 2026 r.}

\end{document}
